\documentclass{meetingminutes}
\meetingcommittee{คณะกรรมการบริหารวิชาการ ประจำคณะเทคโนโลยีอุตสาหกรรม มหาวิทยาลัยราชภัฏอุตรดิตถ์}
\meetingnumber{๒}{๒๕๖๙}
\meetingdate{วันที่ ๒๕ กุมภาพันธ์ พ.ศ. ๒๕๖๙}
\meetingplace{ห้องประชุมคณะเทคโนโลยีอุตสาหกรรม}
\meetingstart{๑๓.๐๐ น.}
\meetingend{๑๕.๐๐ น.}

\begin{document}
\maketitle

\psection{ผู้มาประชุม}
\begin{participants}
  \pmember{1}{ผู้ช่วยศาสตราจารย์ ดร.ชัชพล เกษวิริยะกิจ}{คณบดีคณะเทคโนโลยีอุตสาหกรรม}{ประธาน}
  \pmember{2}{ผู้ช่วยศาสตราจารย์ ดร.พิทักษ์ คล้ายชม}{รองคณบดีฝ่ายวิชาการและวิเทศสัมพันธ์}{กรรมการ}
  \pmember{3}{ผู้ช่วยศาสตราจารย์ ดร.พิศิษฐ์ นาคใจ}{ประธานหลักสูตร วศ.ม. วิศวกรรมคอมพิวเตอร์และปัญญาประดิษฐ์}{กรรมการ}
  \pmember{4}{กรรมการบริหารวิชาการประจำคณะ}{คณะเทคโนโลยีอุตสาหกรรม}{กรรมการ}
\end{participants}

\psection{ผู้ไม่มาประชุม}
\noindent ไม่มี\par

\startmeeting
\openingchair{ผู้ช่วยศาสตราจารย์ ดร.ชัชพล เกษวิริยะกิจ}{กล่าวเปิดการประชุมและดำเนินการประชุมตามระเบียบวาระ ดังนี้}

% ============================================================
\agenda{๑}{เรื่องแจ้งเพื่อทราบ}
% ============================================================

ประธานแจ้งให้ที่ประชุมทราบว่า หลักสูตรวิศวกรรมศาสตรมหาบัณฑิต สาขาวิชาวิศวกรรมคอมพิวเตอร์และปัญญาประดิษฐ์
เสนอขอปรับปรุงแก้ไขหลักสูตรตาม สมอ. ๐๘ โดยขอเพิ่มอาจารย์ประจำหลักสูตรจำนวน ๓ คน
ซึ่งผ่านมติเห็นชอบจากคณะกรรมการบริหารหลักสูตรฯ ครั้งที่ ๑/๒๕๖๙ (๗ มกราคม ๒๕๖๙) เรียบร้อยแล้ว

\noindent\textbf{มติที่ประชุม}\quad ที่ประชุมรับทราบ

% ============================================================
\agenda{๒}{พิจารณาให้ความเห็นชอบการปรับปรุงแก้ไขหลักสูตร วศ.ม. วิศวกรรมคอมพิวเตอร์และปัญญาประดิษฐ์
ตาม สมอ. ๐๘}
% ============================================================

ผู้ช่วยศาสตราจารย์ ดร.พิทักษ์ คล้ายชม (รองคณบดีฝ่ายวิชาการ) นำเสนอรายละเอียดการปรับปรุงแก้ไขหลักสูตร
ตาม สมอ. ๐๘ ต่อที่ประชุม ดังนี้

\begin{thailist}
  \item \textbf{เหตุผลในการปรับปรุง:} หลักสูตรต้องการอาจารย์ผู้มีความเชี่ยวชาญด้านคณิตศาสตร์ประยุกต์
  วิศวกรรมไฟฟ้า ระบบควบคุมอัตโนมัติ และเทคโนโลยีสิ่งแวดล้อม เพื่อเสริมสมรรถนะของนักศึกษา
  ให้ครอบคลุมทั้งด้านวิศวกรรมคอมพิวเตอร์และปัญญาประดิษฐ์อย่างสมบูรณ์
  \item \textbf{อาจารย์ที่ขอเพิ่ม:} ผศ.ดร.พิทักษ์ คล้ายชม (ลำดับ ๘), ผศ.ดร.รัฐพล ดุลยะลา (ลำดับ ๙),
  ผศ.ดร.วีระพล คงนุ่น (ลำดับ ๑๐) — ทุกคนมีวุฒิปริญญาเอกและผลงานวิชาการตรงตามเกณฑ์ กมอ. ๒๕๖๕
  \item \textbf{มีผลตั้งแต่:} ภาคเรียนที่ ๑ ปีการศึกษา ๒๕๖๘ เป็นต้นไป
\end{thailist}

ผู้ช่วยศาสตราจารย์ ดร.พิศิษฐ์ นาคใจ ชี้แจงเพิ่มเติมว่า อาจารย์ทั้ง ๓ คน ไม่มีความซ้ำซ้อนกับหลักสูตรอื่น
และได้ผ่านการตรวจสอบความซ้ำซ้อนจากกองบริการการศึกษาเรียบร้อยแล้ว

ที่ประชุมร่วมกันพิจารณาและเห็นว่าการปรับปรุงดังกล่าวมีความเหมาะสม

\resolution{ที่ประชุมมีมติเห็นชอบการปรับปรุงแก้ไขหลักสูตรวิศวกรรมศาสตรมหาบัณฑิต
สาขาวิชาวิศวกรรมคอมพิวเตอร์และปัญญาประดิษฐ์ ตาม สมอ. ๐๘ และให้นำเสนอต่อ
คณะกรรมการสภาวิชาการเพื่อพิจารณาต่อไป}

% ============================================================
\agenda{๓}{เรื่องอื่น ๆ}
% ============================================================

ไม่มีวาระเพิ่มเติม

\noindent\textbf{มติที่ประชุม}\quad ที่ประชุมรับทราบ

\adjournmeeting

\begin{sigblock}
\typist{ผู้ช่วยศาสตราจารย์ ดร.พิทักษ์ คล้ายชม}
\reviewer{ผู้ช่วยศาสตราจารย์ ดร.พิศิษฐ์ นาคใจ}{กรรมการ}
\chairsig{ผู้ช่วยศาสตราจารย์ ดร.ชัชพล เกษวิริยะกิจ}{ประธานที่ประชุม}
\end{sigblock}

\end{document}
